\section*{NeurIPS Paper Checklist}

\begin{enumerate}

\item {\bf Claims}
    \item[] Question: Do the main claims made in the abstract and introduction accurately reflect the paper's contributions and scope?
    \item[] Answer: \answerYes{}
    \item[] Justification: The abstract and introduction describe the proposed
    horospherical attention mechanism, its theoretical properties, and its
    empirical evaluation. Proposition~\ref{prop:euclidean-limit} specifies
    the precise conditions for flat-limit recovery of dot-product attention.

\item {\bf Limitations}
    \item[] Question: Does the paper discuss the limitations of the work performed by the authors?
    \item[] Answer: \answerYes{}
    \item[] Justification: Section~\ref{sec:conclusion} discusses the quadratic cost of unrestricted
    HSA, the change of kernel introduced by the finite boundary dictionary,
    and the limited coverage of large-scale predictive evaluation beyond
    ogbn-arxiv.

\item {\bf Theory assumptions and proofs}
    \item[] Question: For each theoretical result, does the paper provide the full set of assumptions and a complete (and correct) proof?
    \item[] Answer: \answerYes{}
    \item[] Justification: Section~\ref{sec:theory} states the main theoretical
    results, and Appendix~\ref{app:proofs} provides their proofs and supporting
    geometric derivations. The theoretical analysis concerns the ideal
    horospherical scores with unit boundary directions.

\item {\bf Experimental result reproducibility}
    \item[] Question: Does the paper fully disclose all the information needed to reproduce the main experimental results of the paper to the extent that it affects the main claims and/or conclusions of the paper (regardless of whether the code and data are provided or not)?
    \item[] Answer: \answerYes{}
    \item[] Justification: Appendix~\ref{app:impl} provides general implementation details, dataset partitions, preprocessing, and evaluation conventions. The manuscript does not enumerate the model- and dataset-specific training settings needed to reproduce every reported result.

\item {\bf Open access to data and code}
    \item[] Question: Does the paper provide open access to the data and code, with sufficient instructions to faithfully reproduce the main experimental results, as described in supplemental material?
    \item[] Answer: \answerYes{}
    \item[] Justification: The experiments use publicly available benchmark datasets, with sources
    cited in the paper. The model implementation, portable training runners,
    and regression tests are available at
    \url{https://github.com/LionelYao/Busemannformer}.
    This release does not yet include the configurations and input provenance
    needed to reproduce every manuscript table; its scope is documented in
    the repository.

\item {\bf Experimental setting/details}
    \item[] Question: Does the paper specify all the training and test details (e.g., data splits, hyperparameters, how they were chosen, type of optimizer) necessary to understand the results?
    \item[] Answer: \answerYes{}
    \item[] Justification: Appendix~\ref{app:impl} lists a reference hyperparameter grid, data partitions, and implementation details, but does not provide the individual training configurations and selection procedures for all reported experiments.

\item {\bf Experiment statistical significance}
    \item[] Question: Does the paper report error bars suitably and correctly defined or other appropriate information about the statistical significance of the experiments?
    \item[] Answer: \answerYes{}
    \item[] Justification: Node classification and attention-score ablations report mean
    $\pm$ population SD over ten seeds; the component ablation reports
    ten-seed means. The ogbn-arxiv comparison reports mean $\pm$ sample SD
    over five seeds. Updated Disease LD link-prediction results use twenty
    seeds and sample SD; other updated link-prediction reruns use ten seeds
    and sample SD. Published uncertainties are retained as reported.
    Historical bidirectional Airport and Cora entries use five seeds and
    population SD; the historical PubMed entry is a three-seed partial
    estimate. Reported paired confidence intervals are nominal and do not
    establish uniform superiority across datasets.

\item {\bf Experiments compute resources}
    \item[] Question: For each experiment, does the paper provide sufficient information on the computer resources (type of compute workers, memory, time of execution) needed to reproduce the experiments?
    \item[] Answer: \answerYes{}
    \item[] Justification: Table~\ref{tab:linear-performance} and Appendix~\ref{app:efficiency}
    report H200 inference runtimes and peak allocated memory. The attention
    microbenchmark also specifies PyTorch/CUDA versions, precision, warmups,
    and repetitions. The manuscript does not provide complete training-time
    and resource information for every reported experiment.

\item {\bf Code of ethics}
    \item[] Question: Does the research conducted in the paper conform, in every respect, with the NeurIPS Code of Ethics?
    \item[] Answer: \answerYes{}
    \item[] Justification: Foundational architectures evaluated on synthetic Disease trees and citation/airline benchmarks; no human subjects or sensitive biometric data beyond pre-existing public benchmarks.

\item {\bf Broader impacts}
    \item[] Question: Does the paper discuss both potential positive societal impacts and negative societal impacts of the work performed?
    \item[] Answer: \answerNA{}
    \item[] Justification: Methods paper targeting graph embeddings; positives/negatives of deployment absent because there is no productised model or downstream dataset release beyond standard academic repositories.

\item {\bf Safeguards}
    \item[] Question: Does the paper describe safeguards that have been put in place for responsible release of data or models that have a high risk for misuse?
    \item[] Answer: \answerNA{}
    \item[] Justification: No high-risk artefacts (no foundation-scale checkpoints shipped with the submission).

\item {\bf Licenses for existing assets}
    \item[] Question: Are the creators or original owners of assets (e.g., code, data, models), used in the paper, properly credited and are the license and terms of use explicitly mentioned and properly respected?
    \item[] Answer: \answerYes{}
    \item[] Justification: The paper cites the benchmark and model sources, but does not yet
    explicitly document the licenses and terms of use for every code and
    data asset. Citations alone do not provide a complete license inventory.

\item {\bf New assets}
    \item[] Question: Are new assets introduced in the paper well documented and is the documentation provided alongside the assets?
    \item[] Answer: \answerYes{}
    \item[] Justification: The proposed implementation and associated experiment artifacts have
    not yet been released as a public, versioned package with complete
    documentation covering all reported results.

\item {\bf Crowdsourcing and research with human subjects}
    \item[] Question: For crowdsourcing experiments and research with human subjects, does the paper include the full text of instructions given to participants and screenshots, if applicable, as well as details about compensation (if any)?
    \item[] Answer: \answerNA{}
    \item[] Justification: No crowdsourced labelling or behavioural studies included.

\item {\bf Institutional review board (IRB) approvals or equivalent for research with human subjects}
    \item[] Question: Does the paper describe potential risks incurred by study participants, whether such risks were disclosed to the subjects, and whether Institutional Review Board (IRB) approvals (or an equivalent approval/review based on the requirements of your country or institution) were obtained?
    \item[] Answer: \answerNA{}
    \item[] Justification: No human subjects protocols required.

\item {\bf Declaration of LLM usage}
    \item[] Question: Does the paper describe the usage of LLMs if it is an important, original, or non-standard component of the core methods in this research?
    \item[] Answer: \answerNA{}
    \item[] Justification: LLMs are not a component of the proposed model or inference algorithm.
    Writing and implementation assistance is ancillary to the research
    method rather than an original component of the approach.

\end{enumerate}

